\section{Conclusion}
\label{sec:concl}

We have argued that incomplete histories should be treated as a prediction condition rather than as data to be repaired, and we have supported the argument with a protocol and a model. TrajGap-Bench freezes a single Mixed contract on the public TrajImpute release---one split construction, one checkpoint rule, two metrics, five diagnostic stratifications---so that methods which consume incomplete histories can be compared on identical terms. GAST carries an observation-evidence state through frame-wise relational attention, a reliability-modulated causal state-space encoder and an anchor-conditioned goal decoder, and it never imputes a coordinate.

On 302{,}886 test samples from five scenes, GAST obtains the best macro-averaged best-of-20 error among the compared methods (0.3168/0.4831~m), 9.4\%/7.6\% below the strongest external baseline. Its advantage over the strongest internal baseline is small in aggregate (1.7\%/1.8\%) but structured: it grows monotonically with missing-frame count up to six frames, is largest on terminal gaps, and reverses when a single observation survives. We have also reported two results that qualify the contribution: the improvement is not uniform across scenes, and it is accompanied by an 11--13\% degradation of top-1 ranking, which we trace to the decoder and leave as an open problem.

The next steps follow directly from the limitations. A multi-seed study with trajectory-clustered statistics is needed to establish whether the margin over M0 is stable. A four-arm ablation should separate the encoder and decoder contributions. A saturation rule for evidence age should address the reversal at the oldest anchor. And the same protocol should be extended to other observation geometries and to trajectories produced by a real perception stack, where the missingness distribution is not synthetic.
